% GENERATED FILE. Edit canonical lessons, then run npm run generate:books.
% canonical-source-hash: fnv1a64:adb1a699b0e8291b
% canonical-lessons: PT-C481-sobrancelha, PT-C481-pestana, PT-C481-palpebra, PT-C481-narina, PT-C481-gengiva

\chapter{Olhos e boca}
\label{ch:pt-olhos-e-boca}
\hlchaptermodality{481}

\begin{chapteropening}
\textbf{By the end of this chapter:} \emph{I can name the small parts of the face.}

Complete the last lesson of chapter 481: I can name the small parts of the face.
\end{chapteropening}

\section[a sobrancelha]{a sobrancelha — eyebrow}

\label{lesson:PT-C481-sobrancelha}



\begin{quote}
Before the new one: say the Portuguese for to embarrass, then the Portuguese for ashamed.
\end{quote}

\subsection*{You'll want to know: a sobrancelha}
\textbf{a sobrancelha} — \textquotedblleft{}eyebrow\textquotedblright{}.

\begin{quote}
Quando fica surpreendido, levanta as sobrancelhas.
\end{quote}

\begin{quote}
\emph{When he is surprised, he raises his eyebrows.}
\end{quote}

\emph{Franzir as sobrancelhas} — to frown.

\begin{cousinweb}
\emph{sobrancelha} is inherited, through Old Portuguese \emph{sobrencella}, from Latin \emph{supercilia} \textquotedblleft{}eyebrows\textquotedblright{} (EtymDB), built on \emph{super} \textquotedblleft{}above\textquotedblright{} and \emph{cilium} \textquotedblleft{}eyelid\textquotedblright{}.
\end{cousinweb}

\begin{grammarlens}[title={What you've built}]
One more word for telling, explaining and describing at B1.
\end{grammarlens}

\subsection*{Your turn}
\begin{itemize}
\raggedright
  \item \emph{Say it:} \emph{a sobrancelha}
  \item \emph{Say it:} \emph{a sobrancelha}, once more
  \item \emph{Say it:} a sentence of your own with \emph{a sobrancelha}
  \item \emph{Recall:} say the Portuguese for to embarrass, then the Portuguese for ashamed, then say \emph{a sobrancelha} again
\end{itemize}

\begin{tcolorbox}[breakable,colback=teal!4,colframe=teal!35!black,title={Before you move on}]
What does a sobrancelha mean? (\textquotedblleft{}Eyebrow\textquotedblright{}.) Say it once more.
\end{tcolorbox}
\section[a pestana]{a pestana — eyelash}

\label{lesson:PT-C481-pestana}



\begin{quote}
Before the new one: say the Portuguese for ashamed, then the Portuguese for eyebrow.
\end{quote}

\subsection*{You'll want to know: a pestana}
\textbf{a pestana} — \textquotedblleft{}eyelash\textquotedblright{}.

\begin{quote}
O bebé tem pestanas muito compridas.
\end{quote}

\begin{quote}
\emph{The baby has very long eyelashes.}
\end{quote}

In Brazil: \emph{o cílio}. \emph{Queimar as pestanas} — to study late into the night.

\begin{cousinweb}
\emph{pestana} is already found in Old Portuguese (EtymDB), but its origin is not settled.
\end{cousinweb}

\begin{grammarlens}[title={What you've built}]
One more word for telling, explaining and describing at B1.
\end{grammarlens}

\subsection*{Your turn}
\begin{itemize}
\raggedright
  \item \emph{Say it:} \emph{a pestana}
  \item \emph{Say it:} \emph{a pestana}, once more
  \item \emph{Say it:} a sentence of your own with \emph{a pestana}
  \item \emph{Recall:} say the Portuguese for ashamed, then the Portuguese for eyebrow, then say \emph{a pestana} again
\end{itemize}

\begin{tcolorbox}[breakable,colback=teal!4,colframe=teal!35!black,title={Before you move on}]
What does a pestana mean? (\textquotedblleft{}Eyelash\textquotedblright{}.) Say it once more.
\end{tcolorbox}
\section[a pálpebra]{a pálpebra — eyelid}

\label{lesson:PT-C481-palpebra}



\begin{quote}
Before the new one: say the Portuguese for eyebrow, then the Portuguese for eyelash.
\end{quote}

\subsection*{You'll want to know: a pálpebra}
\textbf{a pálpebra} — \textquotedblleft{}eyelid\textquotedblright{}.

\begin{quote}
Com sono, as pálpebras ficam pesadas.
\end{quote}

\begin{quote}
\emph{When you are sleepy, your eyelids get heavy.}
\end{quote}

A learned word; in everyday speech people talk about \emph{os olhos}.

\begin{cousinweb}
\emph{pálpebra} is a learned borrowing from Latin \emph{palpebra} \textquotedblleft{}eyelid\textquotedblright{} (EtymDB), related to \emph{palpāre} \textquotedblleft{}to stroke gently\textquotedblright{}.
\end{cousinweb}

\begin{grammarlens}[title={What you've built}]
One more word for telling, explaining and describing at B1.
\end{grammarlens}

\subsection*{Your turn}
\begin{itemize}
\raggedright
  \item \emph{Say it:} \emph{a pálpebra}
  \item \emph{Say it:} \emph{a pálpebra}, once more
  \item \emph{Say it:} a sentence of your own with \emph{a pálpebra}
  \item \emph{Recall:} say the Portuguese for eyebrow, then the Portuguese for eyelash, then say \emph{a pálpebra} again
\end{itemize}

\begin{tcolorbox}[breakable,colback=teal!4,colframe=teal!35!black,title={Before you move on}]
What does a pálpebra mean? (\textquotedblleft{}Eyelid\textquotedblright{}.) Say it once more.
\end{tcolorbox}
\section[a narina]{a narina — nostril}

\label{lesson:PT-C481-narina}



\begin{quote}
Before the new one: say the Portuguese for eyelash, then the Portuguese for eyelid.
\end{quote}

\subsection*{You'll want to know: a narina}
\textbf{a narina} — \textquotedblleft{}nostril\textquotedblright{}.

\begin{quote}
O ar entra pelas narinas.
\end{quote}

\begin{quote}
\emph{Air goes in through the nostrils.}
\end{quote}

Usually plural: \emph{as narinas}. The whole nose is \emph{o nariz}.

\begin{cousinweb}
\emph{narina} is inherited from Vulgar Latin \emph{narīna} (EtymDB), built on Latin \emph{nāris} \textquotedblleft{}a nostril\textquotedblright{} (Lewis \& Short).
\end{cousinweb}

\begin{grammarlens}[title={What you've built}]
One more word for telling, explaining and describing at B1.
\end{grammarlens}

\subsection*{Your turn}
\begin{itemize}
\raggedright
  \item \emph{Say it:} \emph{a narina}
  \item \emph{Say it:} \emph{a narina}, once more
  \item \emph{Say it:} a sentence of your own with \emph{a narina}
  \item \emph{Recall:} say the Portuguese for eyelash, then the Portuguese for eyelid, then say \emph{a narina} again
\end{itemize}

\begin{tcolorbox}[breakable,colback=teal!4,colframe=teal!35!black,title={Before you move on}]
What does a narina mean? (\textquotedblleft{}Nostril\textquotedblright{}.) Say it once more.
\end{tcolorbox}
\section[a gengiva]{a gengiva — gum (in the mouth)}

\label{lesson:PT-C481-gengiva}



\begin{quote}
Before the new one: say the Portuguese for eyelid, then the Portuguese for nostril.
\end{quote}

\subsection*{You'll want to know: a gengiva}
\textbf{a gengiva} — \textquotedblleft{}gum (in the mouth)\textquotedblright{}.

\begin{quote}
Escovo os dentes com cuidado para não magoar as gengivas.
\end{quote}

\begin{quote}
\emph{I brush my teeth carefully so as not to hurt my gums.}
\end{quote}

Chewing gum is a different word: \emph{a pastilha elástica}.

\begin{cousinweb}
\emph{gengiva} is inherited, through Old Portuguese, from Latin \emph{gingīva} \textquotedblleft{}the gums\textquotedblright{} (EtymDB).
\end{cousinweb}

\begin{grammarlens}[title={What you've built}]
One more word for telling, explaining and describing at B1.
\end{grammarlens}

\subsection*{Your turn}
\begin{itemize}
\raggedright
  \item \emph{Say it:} \emph{a gengiva}
  \item \emph{Say it:} \emph{a gengiva}, once more
  \item \emph{Say it:} a sentence of your own with \emph{a gengiva}
  \item \emph{Recall:} say the Portuguese for eyelid, then the Portuguese for nostril, then say \emph{a gengiva} again
\end{itemize}

\begin{tcolorbox}[breakable,colback=teal!4,colframe=teal!35!black,title={Before you move on}]
What does a gengiva mean? (\textquotedblleft{}Gum (in the mouth)\textquotedblright{}.) Say it once more.
\end{tcolorbox}
